
\documentclass[12pt]{article}

\usepackage[margin=1in]{geometry}
\usepackage{longtable}
\usepackage{hyperref}
\usepackage{enumitem}
\usepackage{array}

\title{Scientific Forgery Image Detection\\SRS Version 3}
\author{Janelle Rivera \and Eric Marroquin \and Alexis Flores \and Alex Lam}
\date{May 8, 2026}

\begin{document}

\maketitle

\tableofcontents
\newpage

\section{Introduction}

\subsection{Purpose}

This document defines the requirements for the Scientific Forgery Image Detection
System. The system is designed to detect copy-move forgery in biomedical and
scientific images, where regions of an image are duplicated to manipulate research
results.

The system will allow users to upload scientific images through a web interface. These
images will be processed by a backend system designed to support future detection
functionality.

\subsection{Intended Audience}

This document is intended for:

\begin{itemize}
    \item Software developers working on the system
    \item Course instructors evaluating the project
    \item Students involved in system design
    \item Future maintainers of the software
\end{itemize}

\subsection{Overall Description}

The system is a web-based application consisting of:

\begin{itemize}
    \item A frontend interface for image uploads and displaying results
    \item A backend API for handling image processing requests
    \item A detection system capable of identifying potentially duplicated regions in biomedical images
\end{itemize}

At this stage (Snapshot 4), the system includes:

\begin{itemize}
    \item Frontend image upload functionality
    \item Backend API integration
    \item Initial fraudulent region detection workflow
    \item Suspicious region result display
    \item Final interface and workflow improvements
\end{itemize}

\section{External Interface Requirements}

\subsection{User Interface Requirements}

The system provides a web interface with the following components:

\begin{itemize}
    \item Homepage with project description
    \item Image upload page
    \item File selection input for biomedical images
    \item Submit button for analysis requests
    \item Image preview functionality
    \item Detection results display section
    \item Suspicious region highlighting output
\end{itemize}

\subsection{Software Interface Requirements}

\textbf{Frontend:}

\begin{itemize}
    \item React.js / HTML
    \item HTML/CSS
    \item JavaScript
\end{itemize}

\textbf{Backend:}

\begin{itemize}
    \item Python
    \item FastAPI
    \item Uvicorn
    \item OpenCV
    \item NumPy
    \item Pillow
\end{itemize}

\textbf{Communication:}

\begin{itemize}
    \item HTTP requests between frontend and backend
    \item JSON-based API responses
    \item Base64 encoded result images
\end{itemize}

\subsection{Hardware Interface Requirements}

No specialized hardware is required. The system runs on standard computing devices capable
of running a web browser and Docker containers.

\subsection{Communication Requirements}

\begin{itemize}
    \item REST API communication between frontend and backend
    \item Biomedical image upload through HTTP requests
    \item Detection analysis responses returned in JSON format
    \item Frontend display of suspicious region analysis results
\end{itemize}

\section{Legal and Ethical Considerations}

\subsection{Data Integrity}

The system processes scientific images, which may contain sensitive or
research-related content. All uploaded images are handled securely and only used for
analysis purposes.

\subsection{Privacy Considerations}

\begin{itemize}
    \item Uploaded images are used only for detection purposes
    \item No user personal data is required
    \item Uploaded images are deleted from the server after analysis
    \item Data should not be shared externally without permission
\end{itemize}

\subsection{Ethical Considerations}

\begin{itemize}
    \item The system is intended to support scientific integrity and image analysis research
    \item Detection results should not be considered definitive proof of misconduct
    \item False positives or false negatives may occur due to limitations in early-stage detection methods
    \item Human review is still necessary when evaluating scientific images
\end{itemize}

\section{Glossary}

\begin{longtable}{|p{4cm}|p{9cm}|}
\hline
\textbf{Term} & \textbf{Definition} \\
\hline
API & Interface used for communication between software components \\
\hline
Backend & Server-side system that processes requests \\
\hline
Frontend & User interface of the application \\
\hline
Copy-Move Forgery & Image manipulation technique where parts of an image are duplicated \\
\hline
FastAPI & Python framework for building APIs \\
\hline
React.js & JavaScript library for building user interfaces \\
\hline
HTTP & Protocol used for web communication \\
\hline
JSON & Lightweight data format used for API responses \\
\hline
OpenCV & Open-source computer vision library used for image analysis \\
\hline
ORB & Oriented FAST and Rotated BRIEF --- keypoint detection algorithm \\
\hline
Keypoint & A distinctive location in an image used for matching \\
\hline
\end{longtable}

\section{Versions Table}

\begin{longtable}{|p{1.5cm}|p{2.5cm}|p{5.5cm}|p{4cm}|}
\hline
\textbf{Version} & \textbf{Date} & \textbf{Description} & \textbf{Authors} \\
\hline
1.0 & May 8, 2026 & Initial SRS created for Snapshot 1 & Janelle Rivera, Eric Marroquin, Alexis Flores, Alex Lam \\
\hline
1.1 & May 10, 2026 & Updated frontend implementation progress and revised interface requirements & Janelle Rivera, Eric Marroquin, Alexis Flores, Alex Lam \\
\hline
1.2 & May 12, 2026 & Added fraudulent region detection requirements and backend/frontend integration updates & Janelle Rivera, Eric Marroquin, Alexis Flores, Alex Lam \\
\hline
\end{longtable}

\section{References}

\begin{itemize}
    \item FastAPI Documentation: \url{https://fastapi.tiangolo.com/}
    \item React Documentation: \url{https://react.dev/}
    \item OpenCV Documentation: \url{https://opencv.org/}
    \item OWASP API Security Guidelines: \url{https://owasp.org/www-project-api-security/}
\end{itemize}

\end{document}
